\documentclass[11pt,aspectratio=169]{ctexbeamer}

\usetheme{default}
\setbeamertemplate{navigation symbols}{}

\usepackage{amsmath}
\usepackage{booktabs}

% Projector-safe palette: high contrast on common lecture-room projectors.
\definecolor{DeepBlue}{RGB}{11,58,102}
\definecolor{SoftInk}{RGB}{31,41,51}
\definecolor{PaperBg}{RGB}{255,255,255}
\definecolor{PanelBlue}{RGB}{226,237,247}

\setbeamercolor{normal text}{fg=SoftInk,bg=PaperBg}
\setbeamercolor{structure}{fg=DeepBlue}
\setbeamercolor{frametitle}{fg=DeepBlue,bg=PaperBg}
\setbeamercolor{title}{fg=DeepBlue}
\setbeamercolor{block title}{fg=DeepBlue,bg=PanelBlue}
\setbeamercolor{block body}{fg=SoftInk,bg=white}

\newcommand{\Core}[1]{
  \begin{block}{核心结论}
    #1
  \end{block}
}

\newcommand{\LectureMeta}[2]{
  \title{大模型推理系统与实践}
  \subtitle{#1}
  \author{张书豪}
  \institute{华中科技大学计算机学院}
  \date{2027 · #2}
}

\newcommand{\PaperRef}[3]{
  \item \textbf{#1}（#2）\hfill\href{#3}{原文链接}
}

\newcommand{\ResearchSources}[1]{
  \begin{frame}{本讲参考资料}
    \normalsize
    \begin{itemize}
      #1
    \end{itemize}
    \vspace{0.3em}
    \textcolor{SoftInk}{课件中的论文数字均指作者在特定模型、硬件和工作负载下报告的结果，不代表无条件性能。}
  \end{frame}
}

\newcommand{\InternalCaseNote}{
  \small\textcolor{SoftInk}{说明：本页为课程研究案例或教学假设，不代表已发表论文结论；只有在实验包给出模型、硬件、工作负载和原始记录时，数字才可作为结果引用。}
}

% 14 讲细分版的统一参考资料入口。论文数字的解释边界见课程内容基线。

\newcommand{\LectureOneSources}{\ResearchSources{
  \PaperRef{vLLM: Easy, Fast, and Cheap LLM Serving with PagedAttention}{SOSP 2023}{https://doi.org/10.1145/3600006.3613165}
  \PaperRef{SGLang: Efficient Execution of Structured Language Model Programs}{arXiv:2312.07104}{https://arxiv.org/abs/2312.07104}
  \PaperRef{nano-vllm-hust 教学代码}{课程仓库}{https://github.com/intellistream/intro-to-llm-inference-engines}
}}

\newcommand{\LectureTwoSources}{\ResearchSources{
  \PaperRef{vLLM / PagedAttention}{SOSP 2023}{https://doi.org/10.1145/3600006.3613165}
  \PaperRef{Sarathi-Serve}{OSDI 2024}{https://www.usenix.org/conference/osdi24/presentation/agrawal}
  \PaperRef{DistServe}{OSDI 2024}{https://www.usenix.org/conference/osdi24/presentation/zhong-yinmin}
}}

\newcommand{\LectureThreeSources}{\ResearchSources{
  \PaperRef{vLLM / PagedAttention}{SOSP 2023}{https://doi.org/10.1145/3600006.3613165}
  \PaperRef{DistServe}{OSDI 2024}{https://www.usenix.org/conference/osdi24/presentation/zhong-yinmin}
  \PaperRef{Llumnix}{OSDI 2024}{https://www.usenix.org/conference/osdi24/presentation/sun-biao}
}}

\newcommand{\LectureFourSources}{\ResearchSources{
  \PaperRef{Orca}{OSDI 2022}{https://www.usenix.org/conference/osdi22/presentation/yu}
  \PaperRef{Sarathi-Serve}{OSDI 2024}{https://www.usenix.org/conference/osdi24/presentation/agrawal}
  \PaperRef{FastServe}{NSDI 2026}{https://www.usenix.org/conference/nsdi26/presentation/wu-bingyang}
}}

\newcommand{\LectureFiveSources}{\ResearchSources{
  \PaperRef{vLLM / PagedAttention}{SOSP 2023}{https://doi.org/10.1145/3600006.3613165}
  \PaperRef{vAttention}{arXiv:2405.04437}{https://arxiv.org/abs/2405.04437}
  \PaperRef{Rethinking KV Cache Compression}{MLSys 2025}{https://proceedings.mlsys.org/paper_files/paper/2025/hash/26289c647c6828e862e271ca3c490486-Abstract-Conference.html}
  \PaperRef{Mooncake}{FAST 2025}{https://www.usenix.org/conference/fast25/presentation/qin}
}}

\newcommand{\LectureSixSources}{\ResearchSources{
  \PaperRef{DistServe}{OSDI 2024}{https://www.usenix.org/conference/osdi24/presentation/zhong-yinmin}
  \PaperRef{RAGCache}{arXiv:2404.12457}{https://arxiv.org/abs/2404.12457}
  \PaperRef{Mooncake}{FAST 2025}{https://www.usenix.org/conference/fast25/presentation/qin}
  \PaperRef{NEO}{MLSys 2025}{https://proceedings.mlsys.org/paper_files/paper/2025/hash/66a026c0d17040889b50f0dfa650e5e0-Abstract-Conference.html}
}}

\newcommand{\LectureSevenSources}{\ResearchSources{
  \PaperRef{vLLM / PagedAttention}{SOSP 2023}{https://doi.org/10.1145/3600006.3613165}
  \PaperRef{SGLang}{arXiv:2312.07104}{https://arxiv.org/abs/2312.07104}
  \PaperRef{DistServe}{OSDI 2024}{https://www.usenix.org/conference/osdi24/presentation/zhong-yinmin}
  \PaperRef{Mooncake}{FAST 2025}{https://www.usenix.org/conference/fast25/presentation/qin}
}}

\newcommand{\LectureEightSources}{\ResearchSources{
  \PaperRef{FlashAttention}{NeurIPS 2022}{https://arxiv.org/abs/2205.14135}
  \PaperRef{FlexGen}{ICML 2023}{https://proceedings.mlr.press/v202/sheng23a.html}
  \PaperRef{PowerInfer}{SOSP 2024 / arXiv:2312.12456}{https://arxiv.org/abs/2312.12456}
  \PaperRef{FlashInfer}{MLSys 2025}{https://proceedings.mlsys.org/paper_files/paper/2025/hash/dbf02b21d77409a2db30e56866a8ab3a-Abstract-Conference.html}
}}

\newcommand{\LectureNineSources}{\ResearchSources{
  \PaperRef{DEEPSERVE}{USENIX ATC 2025}{https://www.usenix.org/conference/atc25/presentation/hu-junhao}
  \PaperRef{QServe}{MLSys 2025}{https://proceedings.mlsys.org/paper_files/paper/2025/hash/fbe2b2f74a2ece8070d8fb073717bda6-Abstract-Conference.html}
  \PaperRef{vLLM-HUST}{开源维护仓库}{https://github.com/vLLM-HUST/vllm-hust}
}}

\newcommand{\LectureTenSources}{\ResearchSources{
  \PaperRef{vLLM / PagedAttention}{SOSP 2023}{https://doi.org/10.1145/3600006.3613165}
  \PaperRef{Revisiting Disaggregated LLM Serving}{2025 preprint, arXiv:2601.08833}{https://arxiv.org/abs/2601.08833}
  \PaperRef{Rethinking KV Cache Compression}{MLSys 2025}{https://proceedings.mlsys.org/paper_files/paper/2025/hash/26289c647c6828e862e271ca3c490486-Abstract-Conference.html}
}}

\newcommand{\LectureElevenSources}{\ResearchSources{
  \PaperRef{MLPerf Inference}{MLCommons benchmark}{https://mlcommons.org/benchmarks/inference-datacenter/}
  \PaperRef{ACM Artifact Review and Badging}{ACM policy}{https://www.acm.org/publications/policies/artifact-review-and-badging-current}
  \PaperRef{Revisiting Disaggregated LLM Serving}{2025 preprint, arXiv:2601.08833}{https://arxiv.org/abs/2601.08833}
}}

\newcommand{\LectureTwelveSources}{\ResearchSources{
  \PaperRef{nano-vllm-hust}{课程教学代码}{https://github.com/intellistream/intro-to-llm-inference-engines}
  \PaperRef{vLLM-HUST}{HUST 维护的 vLLM fork}{https://github.com/vLLM-HUST/vllm-hust}
  \PaperRef{上游 vLLM}{开源项目}{https://github.com/vllm-project/vllm}
}}

\newcommand{\LectureThirteenSources}{\ResearchSources{
  \PaperRef{课程内容基线与实验材料}{课程仓库}{https://github.com/intellistream/intro-to-llm-inference-engines}
  \PaperRef{DEEPSERVE}{USENIX ATC 2025}{https://www.usenix.org/conference/atc25/presentation/hu-junhao}
  \PaperRef{GhostServe}{MLSys 2026}{https://proceedings.mlsys.org/paper_files/paper/2026/hash/d99e8e80a6c41e148db686918dd7eab3-Abstract-Conference.html}
}}

\newcommand{\LectureFourteenSources}{\ResearchSources{
  \PaperRef{课程完整材料与后续更新}{课程仓库}{https://github.com/intellistream/intro-to-llm-inference-engines}
  \PaperRef{SpecGen}{MLSys 2026}{https://proceedings.mlsys.org/paper_files/paper/2026/hash/66a026c0d17040889b50f0dfa650e5e0-Abstract-Conference.html}
  \PaperRef{GhostServe}{MLSys 2026}{https://proceedings.mlsys.org/paper_files/paper/2026/hash/d99e8e80a6c41e148db686918dd7eab3-Abstract-Conference.html}
}}



\LectureMeta{第 11 讲 实验方法与可验证性}{研究生课程课件}

\begin{document}

\begin{frame}[plain]
  \titlepage
\end{frame}

\begin{frame}{教学目标}
  \Core{本讲讨论实验设计与验证方法，并说明测量能力本身构成系统研究能力的重要部分。}
  \begin{itemize}
    \item 理解 workload、对照组和指标选择。
    \item 区分端到端指标与局部探针。
    \item 学会识别不可信实验结论。
  \end{itemize}
\end{frame}

\begin{frame}{实验设计的最小要素}
  \Core{一个最小可信实验至少需要 workload、对照组、指标和环境说明四部分。}
  \begin{itemize}
    \item 输入分布和服务场景是什么。
    \item 比较对象和控制变量是什么。
    \item 主要指标与辅助探针分别是什么。
    \item 环境和配置是否可复现。
  \end{itemize}
\end{frame}

\begin{frame}{workload 设计的重要性}
  \Core{workload 决定实验实际测量的问题对象；若输入分布与服务场景不匹配，结论将整体偏移。}
  \begin{itemize}
    \item 长上下文和短上下文应区分。
    \item 在线动态到达和离线批处理不能混用。
    \item RAG 和普通对话的链路结构本身就不同。
  \end{itemize}
\end{frame}

\begin{frame}{端到端指标与局部探针}
  \Core{端到端指标告诉你用户感受到什么，局部探针告诉你系统里发生了什么。}
  \begin{itemize}
    \item TTFT、P99、goodput 属于端到端指标。
    \item batch 大小、KV 占用、GPU 利用率属于局部探针。
    \item 研究中通常两类数据都必须报告。
  \end{itemize}
\end{frame}

\begin{frame}{为什么仅报告端到端指标并不足够}
  \Core{仅报告端到端数字可以描述结果现象，但通常不足以解释其形成原因。}
  \begin{itemize}
    \item P99 上升可能来自队列，也可能来自通信。
    \item throughput 下降可能来自 KV 容量，也可能来自调度保守。
    \item 探针指标帮助建立因果解释链。
  \end{itemize}
\end{frame}

\begin{frame}{常见实验陷阱}
  \Core{很多看起来漂亮的结果，问题并不在数字，而在实验边界没有说清。}
  \begin{itemize}
    \item 只报告最好一次，不报告波动。
    \item 输入分布不合理，无法代表目标场景。
    \item 对照组配置不公平或热身不足。
  \end{itemize}
\end{frame}

\begin{frame}{课程项目中的实验设计}
  \Core{课程项目中的常见问题在于实验计划模糊，从而导致结果缺乏解释力。}
  \begin{itemize}
    \item 先选主指标，再补探针指标。
    \item 先写预期会否定当前假设的结果。
    \item 先定义最小可跑场景，再扩到完整 workload。
  \end{itemize}
\end{frame}

\begin{frame}{实验草案的基本结构}
  \Core{实验开始之前，应先明确整体结构，而非在执行过程中临时补充设计。}
  \begin{itemize}
    \item 主问题和待验证假设。
    \item 实验组、对照组与控制变量。
    \item 主要指标、辅助探针和失败标准。
  \end{itemize}
\end{frame}

\begin{frame}{源码穿插：用 nano-vllm-hust 组织最小 benchmark}
  \Core{\texttt{\detokenize{bench.py}} 已经包含最小实验记录需要的配置、warmup 和结果汇总结构。}
  \begin{itemize}
    \item 它显式设定随机种子、请求数量、输入长度与输出长度分布。
    \item 它先执行一次 warmup，再进入正式测量。
    \item 它最终输出总 token、总时间与 throughput，便于说明最小实验记录格式。
  \end{itemize}
\end{frame}

\begin{frame}{源码穿插：教学版实验任务如何设计}
  \Core{结合 \texttt{\detokenize{student_exercises/README.md}}，可以把实验教学拆成“参考实现阅读 + TODO 补全 + 结果解释”三段。}
  \begin{itemize}
    \item 先在参考实现中理解控制点。
    \item 再在 \texttt{\detokenize{student_exercises/bench.py}} 或相关练习文件中补全任务。
    \item 最后要求学生提交结构化实验记录，而不只是运行截图。
  \end{itemize}
\end{frame}

\begin{frame}{讲解：实验方法课的中心不是跑图，而是写出可否定命题}
  \Core{如果实验只是为了证明自己是对的，那课程项目和论文都很容易滑向结果导向叙事。}
  \begin{itemize}
    \item 好实验首先要写清假设，再写清什么结果会直接否定这个假设。
    \item workload、对照组和探针的设计，本质上都是围绕这条可否定命题服务的。
    \item 因此实验方法课训练的是研究纪律，而不是仅仅训练脚本使用能力。
  \end{itemize}
\end{frame}

\begin{frame}{讲解：端到端指标与局部探针为什么必须成对出现}
  \Core{只报端到端结果，结论容易悬空；只报局部探针，结果又容易失去服务意义。}
  \begin{itemize}
    \item 端到端指标告诉我们用户看到了什么，例如 TTFT、P99、goodput。
    \item 局部探针告诉我们系统里发生了什么，例如 batch 变化、KV 占用、编译命中。
    \item 两者结合起来，才能把“现象”变成“因果解释”。
  \end{itemize}
\end{frame}

\begin{frame}{例子：一个看起来完整但其实不可信的实验（1）}
  \Core{实验方法课必须先教学生识别坏实验，否则工具越熟练，错误可能越系统化。}
\begin{block}{表面上很完整}
    有吞吐图、有延迟表、甚至有消融，但 workload 来源和对照配置没有交代清楚。
  \end{block}
\end{frame}

\begin{frame}{例子：一个看起来完整但其实不可信的实验（2）}
  \Core{实验方法课必须先教学生识别坏实验，否则工具越熟练，错误可能越系统化。}
\begin{block}{真正的问题}
    结果无法回答“为什么变好”以及“在哪些条件下不再成立”。
  \end{block}
\end{frame}

\begin{frame}{例子：一个看起来完整但其实不可信的实验（3）}
  \Core{实验方法课必须先教学生识别坏实验，否则工具越熟练，错误可能越系统化。}
\begin{block}{课程结论}
    一个实验是否可信，不取决于图表数量，而取决于因果链和边界是否被写清。
  \end{block}
\end{frame}

\begin{frame}{例子：端到端数字与探针数字不一致时怎么办（1）}
  \Core{最有研究价值的情况往往不是数字全都一致，而是局部证据与端到端现象出现张力。}
\begin{block}{情况一}
    GPU 利用率更高，但 TTFT 更差，说明等待或同步代价可能在放大。
  \end{block}
\end{frame}

\begin{frame}{例子：端到端数字与探针数字不一致时怎么办（2）}
  \Core{最有研究价值的情况往往不是数字全都一致，而是局部证据与端到端现象出现张力。}
\begin{block}{情况二}
    KV 命中率更高，但 goodput 没提升，说明命中本身未必落在主导瓶颈上。
  \end{block}
\end{frame}

\begin{frame}{例子：端到端数字与探针数字不一致时怎么办（3）}
  \Core{最有研究价值的情况往往不是数字全都一致，而是局部证据与端到端现象出现张力。}
\begin{block}{教学意义}
    探针不是为了给端到端结果背书，而是为了帮助发现矛盾并修正假设。
  \end{block}
\end{frame}

\begin{frame}{代码例子：最小 benchmark 模板应包含什么（1）}
  \Core{一份合格的 benchmark，不只是能跑，还应显式携带 workload 和环境信息。}
\begin{block}{代码片段}
    {\ttfamily\small
    \detokenize{seed(0)}\par
    \detokenize{warmup_requests = 10}\par
    \detokenize{num_requests = 100}\par
    \detokenize{measure(ttft, throughput, p99)}\par
    }
  \end{block}
\end{frame}

\begin{frame}{代码例子：最小 benchmark 模板应包含什么（2）}
  \Core{一份合格的 benchmark，不只是能跑，还应显式携带 workload 和环境信息。}
\begin{block}{课堂应追问}
    哪些配置是影响 workload 定义的，哪些配置只是运行时细节？
  \end{block}
\end{frame}

\begin{frame}{代码例子：最小 benchmark 模板应包含什么（3）}
  \Core{一份合格的 benchmark，不只是能跑，还应显式携带 workload 和环境信息。}
\begin{block}{系统含义}
    benchmark 模板的显式化，本身就在减少“实验只存在于一台机器上”的偶然性。
  \end{block}
\end{frame}

\begin{frame}{代码例子：最小日志汇总脚本（1）}
  \Core{实验课需要让学生看到，图表往往来自对原始日志的二次组织。}
\begin{block}{代码片段}
    {\ttfamily\small
    \detokenize{records = load_jsonl("run.log") }\par
    \detokenize{p99 = percentile([r["latency_ms"] for r in records], 99)}\par
    \detokenize{goodput = sum(r["ok"] for r in records) / total_time}\par
    }
  \end{block}
\end{frame}

\begin{frame}{代码例子：最小日志汇总脚本（2）}
  \Core{实验课需要让学生看到，图表往往来自对原始日志的二次组织。}
\begin{block}{课堂应追问}
    如果日志字段设计不对，后面哪些指标将根本算不出来？
  \end{block}
\end{frame}

\begin{frame}{代码例子：最小日志汇总脚本（3）}
  \Core{实验课需要让学生看到，图表往往来自对原始日志的二次组织。}
\begin{block}{课程结论}
    可验证性不只发生在论文写作阶段，而是从日志字段设计就已经开始了。
  \end{block}
\end{frame}

\begin{frame}{例子：消融实验最容易做假的地方（1）}
  \Core{消融的真正任务是隔离变量，而不是制造更多子图。}
\begin{block}{伪消融一}
    一次改了多个控制点，却只关掉其中一个开关就说完成了消融。
  \end{block}
\end{frame}

\begin{frame}{例子：消融实验最容易做假的地方（2）}
  \Core{消融的真正任务是隔离变量，而不是制造更多子图。}
\begin{block}{伪消融二}
    workload 或参数同时变化，导致数字变化无法归因。
  \end{block}
\end{frame}

\begin{frame}{例子：消融实验最容易做假的地方（3）}
  \Core{消融的真正任务是隔离变量，而不是制造更多子图。}
\begin{block}{课程要求}
    每一项消融都要明确写出“被隔离的变量”和“预期影响的指标”。
  \end{block}
\end{frame}

\begin{frame}{课程研究案例（内部实验）：shared workload}
  \Core{真实 benchmark 的 provenance 包括 workload 来源，也包括 adapter 如何把它变成当前系统可运行的输入。}
  \begin{itemize}
    \item KV-Delta 中，shared workload 元数据里的 prompt length 并不等于 Qwen tokenizer token 数。
    \item 原始请求真实 token 数超过 \texttt{\detokenize{max_model_len}}，导致 completion endpoint 返回 400。
    \item 正确做法是记录 token-budget 裁剪、deterministic subset、submodule commit 和失败 artifact。
  \end{itemize}
  \InternalCaseNote
\end{frame}

\begin{frame}{课程研究案例（内部实验）：保留失败运行}
  \Core{failed / blocked run 不是噪音，而是解释系统边界和后续工作的证据。}
  \begin{itemize}
    \item full 40-request shared case 能通过请求阶段，但 consumer trace/save-load 后段过重。
    \item 因此改用 deterministic first-8 subset 完成 m=2/4/8 canonical，不等于完整 workload 已经解决。
    \item 课程要求：把 full-40 标为 blocked，把 first-8 标为 shared subset，不得混写。
  \end{itemize}
  \InternalCaseNote
\end{frame}

\begin{frame}{课程项目的 claim ledger}
  \Core{每个课程项目都应把结果写成可审计 claim，而不是只贴一张性能图。}
  \begin{itemize}
    \item Claim：当前结论能说到什么程度。
    \item Evidence：artifact path、workload source、baseline、指标和设备。
    \item Boundary：invalid/failed/blocked run、未覆盖 workload、尚未解释的反例。
  \end{itemize}
\end{frame}

\begin{frame}{课堂练习：写一页可信实验卡片（1）}
  \Core{实验方法训练最好用一页纸完成：问题、对照、指标、失败条件。}
\begin{block}{练习一}
    写出你的 workload 定义，至少包含输入长度分布、输出长度分布和到达模式。
  \end{block}
\end{frame}

\begin{frame}{课堂练习：写一页可信实验卡片（2）}
  \Core{实验方法训练最好用一页纸完成：问题、对照、指标、失败条件。}
\begin{block}{练习二}
    写出一个主指标和两个辅助探针，并说明三者的因果关系假设。
  \end{block}
\end{frame}

\begin{frame}{课堂练习：写一页可信实验卡片（3）}
  \Core{实验方法训练最好用一页纸完成：问题、对照、指标、失败条件。}
\begin{block}{练习三}
    用一句话写出：什么结果出现时，你会承认当前机制没有兑现预期收益？
  \end{block}
\end{frame}


\begin{frame}{课堂讨论}
  \Core{实验的目标在于判断系统假设是否成立，而非单纯证明预设观点。}
  \begin{itemize}
    \item 为什么单次最好结果价值很低？
    \item 什么时候局部探针最能帮助解释端到端现象？
    \item 哪些结果一旦出现，就说明当前机制假设应该被放弃？
  \end{itemize}
\end{frame}

\begin{frame}{本讲总结}
  \Core{缺乏规范实验设计时，机制难以转化为可信的系统结论。}
  \begin{itemize}
    \item 验证方法本身就是研究贡献的一部分。
    \item 下一讲将把这些方法落到真实代码与复现实验上。
  \end{itemize}
\end{frame}

\LectureElevenSources

\begin{frame}{对应 Tutorial}
  \Core{本讲对应 Tutorial 10。先独立作答，再使用对应 Tutorial 的教师讨论要点核对。}
  \begin{itemize}
    \item 文件：\texttt{\detokenize{build/tutorials/Tutorial_10_实验方法与验证.pdf}}
    \item 动手运行、日志记录与提交要求统一查看 \texttt{\detokenize{build/experiments/}} 和对应 experiment 源目录下的 \texttt{\detokenize{assignment_spring-2027/}}。
  \end{itemize}
\end{frame}

\end{document}
